跨媒介职业转型研究的起点，是把长期而碎片化的公共曝光记录重新组织为可比较的阶段。既有研究多从职业社会学、媒体研究与品牌管理三个方向讨论个体如何在结构变化中维持位置\cite{Super1957,Jenkins2006,Bourdieu1986}，但很少把“平台切换”与“危机修复”放入同一套分析逻辑中。

\subsection{问题意识}

传统媒体时代的职业进阶往往围绕少量核心入口展开，例如封面资源、影视角色与综艺曝光；平台化时代则把表达频率、互动密度与叙事控制权推到更前台的位置。因此，同样是“转型”，其内部机制已经从“角色迁移”扩展为“角色迁移+平台迁移+关系迁移”的三重过程。

\subsection{阶段划分}

本研究把公开案例划分为四个阶段：起步、扩张、危机与重启。阶段划分并不是为了制造静态标签，而是为了便于比较各阶段的主要平台、角色结构和叙事能力。图 \ref{fig:sysu-framework} 展示了本研究的总体框架。

\begin{figure}[H]
    \centering
    \includegraphics[width=0.80\linewidth]{{figures/demo/framework.pdf}}
    \caption{研究框架与推理链路}
    \label{fig:sysu-framework}
\end{figure}

\subsection{研究意义}

这一主题与仓库中的 NSFC 青年项目示例保持同一研究对象与问题意识，但在 thesis 项目中进一步扩展为更长篇幅的结构化写作，目的是证明博士论文模板可以承载公开案例、理论综述、图表、附录和参考文献等完整写作单元。
